\chapter{Work Plan and Conclusion}
\label{ch:plan}

\section{Work Plan}
\label{sec:plan-work}

The work takes four months (Table~\ref{tab:plan}). Everything before the
registration point prepares the experiments; everything after it runs them as
planned.

\begin{table}[htbp]
\centering
\caption{Work plan.}
\label{tab:plan}
\footnotesize
\begin{tabularx}{\textwidth}{@{}>{\raggedright\arraybackslash}p{0.08\textwidth} >{\raggedright\arraybackslash}X >{\raggedright\arraybackslash}p{0.24\textwidth}@{}}
\toprule
Month & Work & Milestone\\
\midrule
1 & The warm substrate and its calibration; the two-tier loop with shared caches; the Evidence Memory across tiers; \Step\ and the handover note; development set: model pair, variants, compute estimate, all settings fixed & a tiered session runs on every development bug\\
2 & E0 and E1: at most $1{,}200$ cheap and $3{,}000$ strong calls, every candidate audited on both paths; E2 by replay & \textbf{registration}; RQ1 to RQ3 analysed\\
3 & E3: handovers on $50$ bugs; E4 on the other $50$ & RQ4 and RQ5 analysed\\
4 & Writing; release of the bot, as a command-line tool with a CI workflow, and of the data & thesis draft\\
\bottomrule
\end{tabularx}
\end{table}

\section{Risks}
\label{sec:plan-risks}

Table~\ref{tab:risks} lists the main risks and how each is handled.

\begin{table}[H]
\centering
\caption{Main risks and how they are handled.}
\label{tab:risks}
\footnotesize
\begin{tabularx}{\textwidth}{@{}>{\raggedright\arraybackslash}p{0.33\textwidth} >{\raggedright\arraybackslash}X@{}}
\toprule
Risk & Handling\\
\midrule
Escalation saves little & the cost breakdown (C1) and the break-even stand on their own; the result is read as in Chapter~\ref{ch:outcomes}\\
\Step\ only matches a fixed switch & reported as a finding, with the simpler rule recommended\\
The handover shows no effect & cold start is the default and keeps its guarantee; a null answers an open question\\
The benchmark is memorised, so the strong tier repairs everything at once & the model pair is chosen on the development set by a fixed rule, and can move down one step\\
The warm path disagrees with ordinary validation & every candidate is also validated on the cold path, whose verdict is the one of record; a bug that fails the calibration falls back\\
The model service changes & versions pinned, answers cached; E2 needs no model\\
The API bill is too high & the upper bound is known in advance, with a fixed order of cuts (Section~\ref{sec:exp-compute})\\
Audits take too long & caps per bug\\
Flaky or order-dependent tests & detected by repeated runs and excluded\\
Four months is tight & the Transcript handover is cut first, then E4; E0 to E2 answer RQ1 to RQ3\\
\bottomrule
\end{tabularx}
\end{table}

\section{Threats to Validity}
\label{sec:plan-threats}

\paragraph{Internal.} Replay assumes that test results and times do not depend
on order, and that a cold strong session does not depend on the cheap tier.
The first is checked on a sample of audits. The second holds by construction,
since what a model is told about a candidate is fixed by the candidate alone;
it is also checked, because the replay runs the bot itself and fails if a
session asks for an answer that was not recorded. Class times are measured
once per candidate and path, so a charge does not capture how the warmth of
the runtime changes within a session. Model answers are cached, so every
comparison is conditional on one table of draws; the replication is over bugs
and sessions, not over re-sampled answers.

\paragraph{External.} One benchmark, one language, one model family and two
tiers; the faulty methods are given to the model. Defects4J may be
memorised~\citep{ramos2024memorizing}, which inflates the cheap tier's repair
share and with it the saving. Prices change, so the break-even is also
reported as a ratio that can be re-read at other prices. CI systems that run
tests in parallel on many machines are not covered, and neither are bugs
without a failing test.

\paragraph{Construct.} Success means a plausible patch, not a correct one. A
cheap model may produce plausible but wrong patches more often than a strong
one, and the bot returns the first plausible patch it finds. The share of
plausible patches that match the developer's fix is reported per tier as a
partial check; judging correctness in general is out of scope.

\section{Future Work}
\label{sec:plan-future}

\paragraph{A benchmark newer than the models.} The clearest extension is a set
of bugs that the models cannot have seen. It removes the memorisation that
inflates the cheap tier's repair share on Defects4J, and a second benchmark
would show whether the break-even moves as its three inputs predict.

\paragraph{More tiers and local models.} A local model has no token price but
costs time and energy, and smaller is not always
cheaper~\citep{vallecillos2026quantization}. A third tier would test whether
the break-even composes.

\paragraph{A learned trigger on top of the verified one.} If \Step\ ends far
from Hindsight, the gap is the room for a predictor such as those of
Section~\ref{sec:rel-f}, used only where the tests leave the decision open.

\paragraph{Other settings.} Python, using
BugsInPy~\citep{widyasari2020bugsinpy}, where there is no build step and one
validation is cheaper; and repository-level issue
resolution~\citep{jimenez2024swebench}, where a failing test must first be
written before a stall can be verified.

\paragraph{Correctness across tiers.} Whether patches from the cheap tier are
correct as often as those from the strong tier, judged by developers or by
held-out tests.

\section{Conclusion}
\label{sec:plan-conclusion}

In LLM-based repair, each round costs a model call and a test run. Known
techniques make the test run cheap. This proposal turns to the model call. It
builds a tiered repair bot that tries a cheap model first and steps up to a
strong model only when the tests show a stall. A Warm Substrate makes cheap
candidates cheap to check. An Evidence Memory keeps what the tests have shown
across the switch. An Escalation Controller applies the rule \Step\ and writes
a handover note. A Budget Guard stops the session and explains why. With a
cold handover the bot loses no repair that the strong model alone would make,
so what the evaluation measures is its cost: by replay on recorded sessions
against a hindsight bound, and in live runs on Defects4J, with the reading of
every outcome fixed in advance.
